%%%%%%%%%%%%%%%%%%%%%%%%%%%%%%%%%%%%%%%%%%%%%%%%%%%%%%%%%%%%%%%%%%%%%%%%%%%%%%%%%%%%%%%%%
% Autor:        Aguilar Enriquez, Paul Sebastian a.k.a. Penserbjorne
% Fecha:        05/02/2017
% Descripcion:  Plantilla base para actividades o tareas.
%%%%%%%%%%%%%%%%%%%%%%%%%%%%%%%%%%%%%%%%%%%%%%%%%%%%%%%%%%%%%%%%%%%%%%%%%%%%%%%%%%%%%%%%%

\documentclass[a4paper,11pt]{article}                 % Papel tamaño carta, texto de 11pt.

\usepackage[top=2cm, bottom=2cm, left=2.2cm, right=2.2cm]{geometry} % Margenes
\usepackage[T1]{fontenc}                              % Indicamos la codificacion de las fuentes.
\usepackage[utf8x]{inputenc}                          % Definimos la codificacion.
\usepackage{lmodern}                                  % Para poder usar acentos.
\usepackage[spanish]{babel}                           % Usaremos idioma español.
\usepackage{amsmath}                                  % Para formulas matematicas.
\usepackage{graphicx}                                 % Para imagenes.
\usepackage{float}                                    % Para posicionar objetos.
\usepackage{booktabs}                                 % Para formatear tablas.
\usepackage{hyperref}                                 % Para enlaces y referencias.
\usepackage{enumerate} 
%\usepackage{colortbl}
%%%%%%%%%%%%%%%%%%%%%%%%%%%%%%%%%%%%%%%%%%%%%%%%%%%%%%%%%%%%%%%%%%%%%%%%%%%%%%%%%%%%%%%%%

% Los logos tienen posiciones relativas al nombre de la escuela.
% Cada imagen esta desplazada con respecto al texto, en este caso nombre de la univseridad.
% No se necesitan paquetes adicionales, el entorno estandar para imagenes de LaTeX puede hacerlo.
% El truco esta en definir una imagen de tamaño cero, asi no afecta al centrar los titulos.
\def\logoUNAM{%
  \begin{picture}(0,0)\unitlength=1cm
    \put (-3.5,-3) {\includegraphics[width=8em]{images/escudo-unam}}
  \end{picture}
}

\def\logoFI{%
  \begin{picture}(0,0)\unitlength=1cm
    \put (0.5,-3) {\includegraphics[width=8em]{images/escudo-fi}}
  \end{picture}
}

%%%%%%%%%%%%%%%%%%%%%%%%%%%%%%%%%%%%%%%%%%%%%%%%%%%%%%%%%%%%%%%%%%%%%%%%%%%%%%%%%%%%%%%%%

\author{Pérez Navarro Maria Yesica - 414039694}  % Autor de la actividad.
\title{Previo 05: Programación en lenguaje\\ Assembly (Ensamblador) 2}                % Titulo de la actividad.
\date{dd/mm/yyyy}                                           % Fecha de entrega.
\def\universidad{Universidad Nacional Autónoma de México}   % Nombre de la universidad.
\def\facultad{Facultad de Ingeniería}                              % Nombre de la facultdad.
\def\semestre{2018-1}                                     % Semestre lectivo.
\def\materia{Lab. Microcontroladores y Microprocesadores - Grupo 03}               % Nombre de la materia y grupo.
\makeatletter

%%%%%%%%%%%%%%%%%%%%%%%%%%%%%%%%%%%%%%%%%%%%%%%%%%%%%%%%%%%%%%%%%%%%%%%%%%%%%%%%%%%%%%%%%

\begin{document}
  
  % Titulo del documento con logos.
  \begin{center}
    \logoUNAM {\Large \universidad} \logoFI\par
    {\large \facultad}\par
    \semestre\par
    \materia\par
    \@author\par
    \@date\par
    \@title
  \end{center}

  \hrulefill\par

  \pagenumbering{gobble}                              % Oculta el numero de pagina.
%  \tableofcontents                                    % Crea el indice o tabla de contenido.

%%%%%%%%%%%%%%%%%%%%%%%%%%%%%%%%%%%%%%%%%%%%%%%%%%%%%%%%%%%%%%%%%%%%%%%%%%%%%%%%%%%%%%%%%

  %\newpage                                            % Inserta una pagina nueva.
  \pagenumbering{arabic}                              % Muestra el numero de pagina.
  
  \section{Seguridad en la Ejecución.}
  \begin{table}[H]
  	\begin{tabular}{|l|l|l|}
  		\hline
  		 & Peligro o fuente de energía & Riesgo asociado  \\ \hline
  		1 & Manejo de Corriente Alterna &Electrochoque    \\ \hline
  		2 & Manejo de corriente Continua & Daño al equipo \\ \hline
  	\end{tabular}
  	\centering
  \end{table}

\section{Objetivos de aprendizaje.}
\begin{itemize}
	\item El alumno aprenderá a utilizar las instrucciones de operaciones aritméticas y lógicas, así como emplear saltos condicionales y llamados a subrutinas. 
\end{itemize}

\section{Material y equipo.}
 
\begin{itemize}
	\item Calculadora en modo “Programador”.
	\item Tarjeta de desarrollo.
	\item CCS IDE. 
\end{itemize}
  

  
\section{Actividad previa.}                   

\begin{enumerate}[a)]
	\item ¿Cómo se obtiene la representación binaria de un número negativo en el formato de complemento a 2? 
	
	\item  ¿Cuál es la longitud en bytes del resultado de una multiplicación no signada (unsigned multiplication) al multiplicar: a) dos valores de 16 bits cada uno? b) dos valores de 32 bits cada uno? 
	
	\item  ¿Qué hace el salto condicional “BEQ etiqueta1” ? 
	
	\item  ¿Qué hace el salto “BL etiqueta2” ? 
	
	\item  Para cada uno de los dos saltos anteriores, ¿se debe regresar al lugar del salto?, ¿cómo?
	
	\item Examine el comportamiento de los segmentos del código proporcionado. Con ayuda de la calculadora, compruebe el resultado de las operaciones.
	
\end{enumerate}


 \section{Desarrollo.}

 Escriba un programa que funcione como una calculadora de números enteros de hasta 32 bits, con operaciones de Suma, Resta, Multiplicación, AND, OR, XOR. 

 Para elegir la operación a realizar, desde el depurador se escribirá un valor en R0 del 1 al 6 (1:suma, 2:resta, etc.). El programa debe hacer llamados a dos subrutinas: una para carga de operandos y otra para guardar resultado.
 
  Los dos operandos se encuentran definidos en memoria RAM y se deben colocar en R1 y R2. El resultado se escribe también en memoria RAM, usando R3. 
  
  Por tanto se debe hacer acceso de lectura y escritura (LDR, STR). 
  
  Los operandos se pueden modificar en cualquier sobre-escribiendo su valor en memoria RAM. 


  
  \section{Cuestionario}
  
  \begin{enumerate}
  	\item Después de ejecutar la instrucción SUBS, ¿qué signo tiene el resultado cuando las bandera C = 0, y cuando C=1? En cada caso, qué valor debe tener la bandera N? 
  	
  	Para un valor en un Registro, en lenguaje ensamblador.

  	\item ¿Cómo se enciende un bit específico sin afectar a los demás? (2 maneras).

  	\item ¿Cómo se apaga un bit específico sin afectar a los demás? (2 maneras).
  	
  	\item ¿Cómo se invierte el estado de un bit específico?.
  	
  	\item Al emplear instrucción de comparación CMP, ¿Qué nos dice la bandera Z?.
  	
  	\item Se quiere saber si los 4 LSB de un byte son ‘0’, ¿cómo emplearía la instrucción TST y qué condición/bandera se debe cumplir/establecer?.
  	
  	\item Se quiere saber si los 4 MSB de un byte son ‘1’, ¿cómo emplearía la instrucción TEQ y qué condición/bandera se debe cumplir/establecer? 
  	
  	\item Para qué sirve la directiva .retain 
  	  
  \end{enumerate}

\section{Conclusiones.}
  \begin{thebibliography}{}                           % Bibliografia
    \bibitem{ref:cita}                                % Etiqueta con la que se hara la referencia o cita.
      Como citar: \url{http://www.cva.itesm.mx/biblioteca/pagina_con_formato_version_oct/apa.htm} % URL de apoyo para citar.

    \bibitem{ref:web1}
      Autor,
      (Fecha de publicacion),
      Titulo, paginas,
      Fecha de recuperacion,
      Sitio web: \url{http://www.google.com}

    \bibitem{ref:github}
      Repositorio del proyecto \url{https://github.com/penserbjorne}
  \end{thebibliography}

%%%%%%%%%%%%%%%%%%%%%%%%%%%%%%%%%%%%%%%%%%%%%%%%%%%%%%%%%%%%%%%%%%%%%%%%%%%%%%%%%%%%%%%%%

\end{document}
